\documentclass{article}
\usepackage{amsmath,amssymb}
\usepackage{xurl}
\usepackage{needspace}
\usepackage[hidelinks]{hyperref}
% Shared commands for notes in docs/paper-gaps/.

\DeclareUnicodeCharacter{2080}{\ifmmode{}_{0}\else\textsubscript{0}\fi}
\DeclareUnicodeCharacter{2081}{\ifmmode{}_{1}\else\textsubscript{1}\fi}
\DeclareUnicodeCharacter{2082}{\ifmmode{}_{2}\else\textsubscript{2}\fi}
\DeclareUnicodeCharacter{2099}{\ifmmode{}_{n}\else\textsubscript{n}\fi}

\newcommand{\Lean}{\textsc{Lean}}
\ExplSyntaxOn
\NewDocumentCommand{\leanid}{m}{
  \ifmmode
    \textsf{\detokenize{#1}}
  \else
    \group_begin:
    \urlstyle{sf}
    \tl_set:Nn \l_tmpa_tl {#1}
    \tl_replace_all:Nnn \l_tmpa_tl {\_} {_}
    \exp_args:NV \path \l_tmpa_tl
    \group_end:
  \fi
}
\ExplSyntaxOff
\providecommand{\tr}{\operatorname{tr}}
\newcommand{\tnleanIssueBase}{https://github.com/LionSR/TNLean/issues/}
\newcommand{\tnleanPrBase}{https://github.com/LionSR/TNLean/pull/}
\newcommand{\tnleanDocBase}{https://sirui-lu.com/TNLean/docs/}
\newcommand{\ghissue}[1]{\href{\tnleanIssueBase#1}{GitHub issue \##1}}
\newcommand{\ghpr}[1]{\href{\tnleanPrBase#1}{PR \##1}}
\newcommand{\doclink}[1]{\href{\tnleanDocBase#1.html}{\path{#1}}}

% Dirac notation and the identity operator, as in blueprint/src/macros/common.tex.
% \providecommand keeps note-local definitions authoritative.
\usepackage{dsfont}
\providecommand{\ket}[1]{|#1\rangle}
\providecommand{\bra}[1]{\langle #1|}
\providecommand{\braket}[2]{\langle #1 | #2 \rangle}
\providecommand{\ketbra}[2]{|#1\rangle\!\langle #2|}
\providecommand{\Id}{\mathds{1}}
\providecommand{\one}{\mathds{1}}
\providecommand{\C}{\mathbb{C}}
\providecommand{\MN}[1]{M_{#1}(\C)}
\providecommand{\MD}{M_D(\C)}

% Verdict marker: \gapnote{<kind>}{<status>}. Kind is the policy
% classification (clarification, local-correction, scope-restriction,
% unfaithful, false-source, open-gap); status is open, wip (elimination
% underway), resolved, or historical. Machine-read by the site index;
% prints a verdict line.
\newcommand{\gapnote}[2]{\par\noindent\textbf{Verdict:} #1 (#2).\par}

\title{Injective PEPS Parent Interactions: the Physical Reconstruction Condition}
\author{}
\date{2026-10-02}
\begin{document}
\maketitle
\gapnote{local-correction}{resolved}

\section{The printed construction}
Section~IV.C.1 of \cite{Cirac2021Matrix} constructs an injective PEPS
from virtual entangled pairs by a product of site maps $A$.
The equation labelled \path{eq:4:deformed-parent-1} transports a virtual
parent interaction $h$ by a chosen site left inverse:
\begin{align}
    h'=(A^{-1}\otimes A^{-1})^\dagger
             h(A^{-1}\otimes A^{-1}).
    \label{eq:injective_parent_bare}
\end{align}
The following sentence asserts that its kernel is the regional PEPS
space.\footnote{Source traceability:
    \path{Papers/2011.12127/TN-Review-main.tex}, lines 2031--2037.}
This kernel assertion needs a physical reconstruction condition when
the injective site maps are rectangular.

\section{A nondegenerate example}
Consider two vertices joined by a bond of dimension $D=1$ and take
$A:\mathbb C\longrightarrow\mathbb C^2$, $A(z)=(z,0)$, at each vertex.
The bond space is nonzero and both site maps are injective.
The normalized virtual Bell vector is the scalar $1$, so its complement
projector is $h=0$. Consequently (\ref{eq:injective_parent_bare}) gives
$h'=0$ on $\mathbb C^2\otimes\mathbb C^2$. Its kernel is four-dimensional,
whereas the actual two-site PEPS space is
$\operatorname{span}\{e_0\otimes e_0\}$.
Thus the failure is present for positive bond dimension and an injective,
nonzero tensor; it is not a zero-dimensional convention.\footnote{Results:
    \leanid{TNLean.PEPS.InjectiveParentReconstructionCounterexample.bareKernel_ne_productSiteRange},
    \leanid{TNLean.PEPS.InjectiveParentReconstructionCounterexample.finrank_productSiteRange}, and
    \leanid{TNLean.PEPS.InjectiveParentReconstructionCounterexample.finrank_bareKernel}.}

\Needspace{10\baselineskip}
\section{The corrected interaction}
Let $F:\mathcal V\longrightarrow\mathcal H$ be the product site map
on a finite region, and let $L:\mathcal H\longrightarrow\mathcal V$
be a left inverse, so $LF=I_{\mathcal V}$. For any positive semidefinite
virtual interaction $H$, put
\begin{align}
    K=L^\dagger H L+(I_{\mathcal H}-FL)^\dagger(I_{\mathcal H}-FL).
    \label{eq:injective_parent_reconstruction}
\end{align}
Both summands are positive semidefinite. Their sum annihilates $\psi$
precisely when $HL\psi=0$ and $FL\psi=\psi$. Indeed,
\begin{align*}
    \langle\psi,K\psi\rangle
      =\|H^{1/2}L\psi\|^2+\|(I_{\mathcal H}-FL)\psi\|^2.
\end{align*}
It follows that $\ker K=F(\ker H)$.
In contrast, the bare interaction satisfies
$\ker(L^\dagger HL)=F(\ker H)\oplus\ker L$.
The additional summand in (\ref{eq:injective_parent_reconstruction})
removes exactly the unused physical directions. It vanishes if $F$ is
bijective, and does not alter the quadratic form on $F(\mathcal V)$.
If $\ker H$ is the actual regional virtual bond space, the corrected
kernel is exactly the regional physical PEPS space.\footnote{Results:
    \leanid{TNLean.PEPS.injectiveRegionParentTransport_posSemidef},
    \leanid{TNLean.PEPS.injectiveRegionParentTransport_mulVec_eq_zero_iff},
    \leanid{TNLean.PEPS.ker_injectiveRegionParentTransport}, and
    \leanid{TNLean.PEPS.isRegionParentInteraction_injectiveRegionParentTransport}.}

\section{Scope of the correction}
This corrects the displayed interaction, not the uniqueness conclusion
for genuine parent interactions with the prescribed exact local kernels.
For injective tensors, uniqueness of the ground state of genuine
nearest-neighbour parent interactions follows from those kernels and the
virtual Bell constraints. Unused physical directions must first be excluded. Every vertex must be covered by a parent region;
isolated vertices require their own single-site conditions.
The reconstruction summand acts on the chosen region. Its existence does
not by itself establish a decomposition into smaller local terms, nor
any infinite-volume spectral-gap statement.

\bibliographystyle{alpha}
\bibliography{references}
\end{document}
